\documentclass[UTF8]{ctexart}

%\RequirePackage{amsmath} \RequirePackage{amssymb}
\usepackage{amscd,amsthm,amsfonts,amssymb,amsmath}
\usepackage[colorlinks=true]{hyperref}
%\usepackage{fullpage}
\usepackage[all]{xy}
\usepackage{mathrsfs}
\usepackage{tikz,mathpazo}
%\usetikzlibrary{quotes,angles}
%\usetikzlibrary{calc}
%\usetikzlibrary{decorations.pathreplacing}
\usepackage{bm}
\usepackage{geometry}
\usepackage{xcolor}
\usepackage{eso-pic}
\usepackage{CJK}
\usepackage{arydshln}
\usepackage{mathptmx}
\usepackage[11pt]{moresize}




\newcommand{\BC}{{\mathbb {C}}}

\newcommand{\BN}{{\mathbb {N}}}
\newcommand{\BQ}{{\mathbb {Q}}}
\newcommand{\BR}{{\mathbb {R}}}
\newcommand{\BZ}{{\mathbb {Z}}}
\newcommand{\bv}{{\bm v}}
\newcommand{\bw}{{\bm w}}
\newcommand{\bu}{{\bm u}}
\newcommand{\bx}{{\bm x}}
\newcommand{\by}{{\bm y}}

\newcommand{\bb}{{\bm b}}
\newcommand{\Coker}{{\mathrm{Coker}}}


\newcommand{\End}{{\mathrm{End}}}
\newcommand{\GL}{{\mathrm{GL}}}

\newcommand{\SO}{{\mathrm{GSO}}}
\newcommand{\Hom}{{\mathrm{Hom}}}

\renewcommand{\Im}{{\mathrm{Im}}}

\newcommand{\Isom}{{\mathrm{Isom}}}
\newcommand{\Id}{{\mathrm{Id}}}



\newcommand{\Ker}{N}
\newcommand{\Ran}{R}
\newcommand{\nullspace}{N}
\newcommand{\range}{R}

\newcommand{\rank}{{\mathrm{rank}}}


\newcommand{\SU}{{\mathrm{SU}}}
\newcommand{\Sym}{{\mathrm{Sym}}}

\newcommand{\sub}{{\mathrm{sub}}}

\newcommand{\tr}{{\mathrm{tr}}}

\newcommand{\matrixx}[4]{\begin{pmatrix}
		#1 & #2 \\ #3 & #4
	\end{pmatrix} }
	
	
	
	\newcommand{\wt}{\widetilde}
	\newcommand{\wh}{\widehat}
	\newcommand{\pp}{\frac{\partial\bar\partial}{\pi i}}
	\newcommand{\pair}[1]{\langle {#1} \rangle}
	\newcommand{\wpair}[1]{\left\{{#1}\right\}}
	\newcommand{\intn}[1]{\left( {#1} \right)}
	\newcommand{\norm}[1]{\|{#1}\|}
	\newcommand{\sfrac}[2]{\left( \frac {#1}{#2}\right)}
	\newcommand{\ds}{\displaystyle}
	\newcommand{\ov}{\overline}
	\newcommand{\incl}{\hookrightarrow}
	\newcommand{\lra}{\longrightarrow}
	\newcommand{\lla}{\longleftarrow}
	\newcommand{\ra}{\rightarrow}
	\newcommand{\imp}{\Longrightarrow}
	\newcommand{\lto}{\longmapsto}
	\newcommand{\bs}{\backslash}
	
	
	
	
	\newtheorem{thm}{定理}
	\newtheorem{cor}[thm]{推论}
	\newtheorem{lem}[thm]{引理}
	\newtheorem{prop}[thm]{命题}
	\newtheorem{defn}[thm]{定义}
	\newtheorem{ques}[thm]{问题}
	\newtheorem{eg}[thm]{例题}
	\newtheorem{ex}[thm]{习题}
	\newtheorem{rmk}[thm]{注释}
		\newtheorem{ntn}[thm]{记号}

	\newcommand{\sk}{\medskip}\newcommand{\bsk}{\bigskip}
	\newcommand{\s}{\sk\noindent}\newcommand{\bigs}{\bsk\noindent}
	
	
	%accented words

	\def\mat(#1,#2,#3,#4){
		\left[\begin{array}{cc}
			#1 & #2 \\ #3 & #4\end{array}
		\right]
	}
		\def\clm(#1,#2){
		\left[\begin{array}{c}
			#1 \\ #2 \end{array}
\right]
	}
		\def\clmn(#1,#2){
		\left[\begin{array}{c}
			#1 \\ \vdots\\ #2 \end{array}
\right]
	}
		\def\clmt(#1,#2,#3){
		\left[\begin{array}{c}
			#1 \\ #2  \\ #3\end{array}
\right]
	}
		\def\clmf(#1,#2,#3){
		\left[\begin{array}{c}
			#1 \\ #2  \\ \vdots \\ #3\end{array}
\right]
	}


		
\begin{document}


\title{习题课材料（五）}

\date{}

\maketitle

\vspace{-1cm}

\noindent{\Large\textbf{注: 带 $\heartsuit $号的习题有一定的难度、比较耗时, 请量力为之.}}

\medskip

\noindent{\bfseries 记号}: 如不加说明，我们只考虑实矩阵。对于矩阵$A$, 它的四个基本子空间是列空间$C(A)$, 零空间$N(A)$, 行空间$C(A^T)$和$A^T$的零空间$N(A^T)$。

\medskip




\begin{ex}
请构造以下满足条件的矩阵，或者指明为什么不存在满足条件的矩阵：
\begin{enumerate}
\item 列空间包含向量$\begin{bmatrix}1\\2\\-3\end{bmatrix},\begin{bmatrix}2\\-3\\5\end{bmatrix}$，零空间包含向量$\begin{bmatrix}1\\1\\1\end{bmatrix}$；
\item 行空间包含向量$\begin{bmatrix}1\\2\\-3\end{bmatrix},\begin{bmatrix}2\\-3\\5\end{bmatrix}$，零空间包含向量$\begin{bmatrix}1\\1\\1\end{bmatrix}$；
\item $A\bm x=\begin{bmatrix}1\\1\\1\end{bmatrix}$有解，且$A$的左零空间包含$\begin{bmatrix}1\\0\\0\end{bmatrix}$；
\item $A$不是零矩阵，且$A$的所有行垂直于$A$的所有列；
%\item $A$非零，且它的列空间等于零空间；
\item $A$的所有列向量的和是零向量，但是所有行向量的和是一个分量均为$1$的向量；
\item $A,B$均为非零的正交投影矩阵，且$A+B$仍是一个正交投影矩阵；
\item $A,B$均为正交投影矩阵，然而$A+B$并不是一个正交投影矩阵；
\item $A,B,C$均为非对角的三阶正交投影矩阵，且$A+B+C=I_3$。
\end{enumerate}
\end{ex}



\begin{ex}
请将以下向量$\bm x$分解成在$\Ker(A)$中的部分与在$C(A^{T})$中的部分的和，然后点积验证它们确实垂直。
\begin{enumerate}
\item $\bm x=\begin{bmatrix}2\\0\end{bmatrix},A=\begin{bmatrix}1&-1\\0&0\\0&0\end{bmatrix}$；
\item $\bm x=\begin{bmatrix}1\\0\end{bmatrix},A=\begin{bmatrix}1&1\\1&1\end{bmatrix}$；
\item $\bm x=\begin{bmatrix}5\\5\\9\end{bmatrix},A=\begin{bmatrix}1&2&3\\2&2&4\\2&3&5\end{bmatrix}$；
\end{enumerate}
\end{ex}



\begin{ex}
考虑矩阵$A=\begin{bmatrix}3&6&6\\4&8&8\end{bmatrix}$。
\begin{enumerate}
\item 计算向着$\begin{bmatrix}3&6&6\\4&8&8\end{bmatrix}$的列空间的正交投影矩阵$P_1$；
\item 计算向着$\begin{bmatrix}3&6&6\\4&8&8\end{bmatrix}$的行空间的正交投影矩阵$P_2$；
\item 计算$P_1AP_2$。结果意不意外？为什么会有如此结果？
\end{enumerate}
\end{ex}





\begin{ex}[$\heartsuit$]
（Fredholm二择一定理）
\begin{enumerate}
\item 证明方程组$A\bm x=\bm b$有解当且仅当$\Ker(A^T)$与$\bm b$正交；
\item 证明$\Ker(A^T)$与$\bm b$不正交当且仅当存在$\bm y$，使得$\bm y^T \bm A =0, \bm y^T \bm b =1$；
\item 说明当方程组$A\bm x=\bm b$无解时，方程组的一个伪解（即一个假定的$A\bm x=\bm b$的解）给出标准矛盾等式$0=1$。

\end{enumerate}
\end{ex}



\begin{ex}
考虑$\mathbb{R}^n$中两个子空间$\mathcal{M},\mathcal{N}$。通过把对应的一组基当成列向量，不妨假设可以找到列满秩矩阵$A,B$，使得$C(A)=\mathcal{M}, C(B)=\mathcal{N}$。
\begin{enumerate}
\item 证明$\mathcal{M}\subseteq\mathcal{N}$当且仅当存在矩阵$X$使得$A=BX$；
\item 仿照集合论中补集的性质，我们期待正交补有类似的性质。假设$\mathcal{M}\subseteq\mathcal{N}$，请证明$\mathcal{N}^\perp\subseteq\mathcal{M}^\perp$。（能否从前一小问，看出一个矩阵角度的证明？）
\end{enumerate}
\end{ex}



\begin{ex}[$\heartsuit $]
考虑$\mathbb{R}^n$中两个子空间$\mathcal{M},\mathcal{N}$。通过把对应的一组基当成列向量，不妨假设可以找到列满秩矩阵$A,B$，使得$C(A)=\mathcal{M}, C(B)=\mathcal{N}$。
\begin{enumerate}
\item $\mathcal{M}+\mathcal{N}$是哪个由$A,B$构造出的矩阵的列空间？这意味着$(\mathcal{M}+\mathcal{N})^\perp$是该矩阵的什么空间？
\item $\mathcal{M}^\perp,\mathcal{N}^\perp$分别是哪个由$A,B$构造出的矩阵的零空间？这意味着$\mathcal{M}^\perp\cap\mathcal{N}^\perp$是哪个由$A,B$构造出的矩阵的零空间？
\item 证明子空间版本的De Morgan定律：$(\mathcal{M}+\mathcal{N})^\perp=\mathcal{M}^\perp\cap\mathcal{N}^\perp,(\mathcal{M}\cap\mathcal{N})^\perp=\mathcal{M}^\perp+\mathcal{N}^\perp$。
\item[注记:] 集合论中有De Morgan定律，它说对于两个子集$X,Y$，$X\cap Y$的补集等于$X$的补集并上$Y$的补集，且$X\cup Y$的补集等于$X$的补集交上$Y$的补集。对于子空间来说，正交补也有类似的性质。
\end{enumerate}
\end{ex}




\begin{ex}
设$A$为列满秩矩阵。
\begin{enumerate}
\item 假设$\bm v=A\bm x$，我们希望找到$\bm y$使得$\bm v=AA^T\bm y$。请找到一个矩阵$C$（使用$A$来构造）使得$C\bm x$就是这里所需要的$\bm y$；（提示：$A^T$未必可逆，但是利用可逆矩阵$A^TA$，可以给它找一个右逆。）
\item 证明$C(AA^T)=C(A)$。
%\item 证明$C(AA^T+BB^T)$是$\mathcal{M}+\mathcal{N}$。（提示：将矩阵和写成$XX^T$的形式，再利用上一小问。）
%\item 证明$\Ker(AA^T+BB^T)=\mathcal{M}^\perp\cap\mathcal{N}^\perp=(\mathcal{M}+\mathcal{N})^\perp$；
%\item 证明或举出反例：$AA^T+BB^T$是到$\mathcal{M}+\mathcal{N}$的正交投影矩阵的若干倍。
\end{enumerate}
\end{ex}




\begin{ex}[$\heartsuit $]
任取$\mathbb{R}^m$中的子空间$\mathcal{M}_1,\mathcal{M}_2$，与$\mathbb{R}^n$中的子空间$\mathcal{N}_1,\mathcal{N}_2$，是否一定存在一个矩阵$A$，使得$C(A)=\mathcal{M}_1,\nullspace(A^{T})=\mathcal{M}_2, C(A^{T})=\mathcal{N}_1,\nullspace(A)=\mathcal{N}_2$？如果并不一定存在，请给四个子空间加上尽量少的条件，使得这样的矩阵一定存在。
\end{ex}




\begin{ex}
日常生活中我们常说的两个平面垂直，是指平面的法向量正交。例如假设墙角是原点，那么竖直的墙面和地板是一对彼此垂直的平面。但是注意，它们之间并不是正交子空间的关系，因为它俩的交集包含的非零向量。我们下面来探索平面垂直与子空间正交之间的关系。
\begin{enumerate}
\item 考虑$A=\begin{bmatrix}1&2\\1&3\\1&2\end{bmatrix},B=\begin{bmatrix}5&4\\6&3\\-5&-4\end{bmatrix}$。 请找到这两个矩阵的列空间交集中的一个非零向量。他们的列空间平面的法向量是否垂直？是否是正交子空间？
\item 找到以上矩阵对应的正交投影矩阵$P_A,P_B$，是否有$P_AP_B=P_BP_A$？
\item 假设$\mathbb{R}^3$内一个二维子空间作为平面的单位法向量是$\bm v$，证明到该平面的正交投影是$I-\bm v\bm v^T$
\item 给定$\mathbb{R}^3$中的任意两个二维子空间，假设其正交投影矩阵是$P_1,P_2$，证明$P_1P_2=P_2P_1$当且仅当两个子空间要么作为平面是垂直的（法向量垂直），要么相等。（对比之下，正交子空间的要求是$P_1P_2=P_2P_1=0$）
\end{enumerate}
\end{ex}




\begin{ex}
给定$\bm a\neq \bm 0$，方程$\bm a^{T}\bm x=b$定义了$\mathbb{R}^n$中的一个超平面。给定一个向量$\bm v\in\mathbb{R}^n$，将它看成是空间中的一个点，超平面上距离点$\bm v$最近的点称为$\bm v$到该超平面的正交投影。我们想求点$\bm v$到这个超平面的正交投影。
\begin{enumerate}
\item 我们显然希望$\bm v$沿着一个垂直于该超平面的法向量$\bm n$的方向移动，直到$\bm v$碰到该超平面为止。这个法向量$\bm n$是什么？
\item 设答案为$\bm v+x\bm n$，请求出所需要的$x$；
\item 将$\bm v=\begin{bmatrix}1\\2\\3\end{bmatrix}$看成是空间$\mathbb{R}^3$中的一个点，且$\bm a=\begin{bmatrix}2\\1\\-1\end{bmatrix}$，求点$\bm v$到平面$\bm a^{T}\bm x=2$的正交投影和点到平面的距离。
\end{enumerate}
\end{ex}




\begin{ex}
假设$A=\begin{bmatrix}0&0&2\\0&3&0\\1&0&1\end{bmatrix},\bm b=\begin{bmatrix}2.1\\3\\2\end{bmatrix}$。我们考虑方程组$A\bm x=\bm b$的解。显然，$\bm x_0=\begin{bmatrix}1\\1\\1\end{bmatrix}$是一个非常接近正确答案的向量。
\begin{enumerate}
\item 注意，$A\bm x=\bm b$有三行，因此等于是在求解三个平面的交集。请将$\bm x_0$向第一个平面$2x_3=2.1$进行投影，得到$\bm x_1$；
\item 请将$\bm x_1$向第二个平面$3x_2=3$进行投影，得到$\bm x_2$；
\item 请将$\bm x_2$向第三个平面$x_1+x_3=2$进行投影，得到$\bm x_3$；
\item 请将$\bm x_3$再向第一个平面$2x_3=2.1$进行投影，得到$\bm x_4$；
\item （$\heartsuit $）在本题这样的情形中，这样无限循环下去，是否会越来越接近正确答案？
\item[注记：] 实践中，我们想要解的方程组$A\bm x=\bm b$可能十分巨大，比方说做CT扫描时，将人体看成1000个像素块（未知数），用10000条射线以不同的方式穿过（不同的线性方程），那么这等价于在$A$为$10000\times 1000$的时候求解线性方程组。可以想象，用Gauss消元法解决需要花费很久很久的时间，并不现实。所以我们需要利用$A$的特殊性质来想办法加速计算。

有时$A$确实有些很有趣的性质。注意，CT扫描时，$A$每一行对应着一条射线，每一列对应着一个待检测的像素块。然而，每条射线一般仅能穿过很极其有限的像素块，而完全不会碰到大多数像素块。因此$A$的每一行存在着大量的零。一个绝大多数元素都是零的矩阵，我们也称为稀疏矩阵。此时，是否有更快速地解方程的方法呢？我们这里介绍的方法叫做Kaczmarz方法。由于$A$是一个稀疏矩阵，每一行存在大量的零。因此投影计算非常快。所以实践中，可以采取这种Kaczmarz方法进行反复投影，直到足够接近正确答案为止。
\end{enumerate}
\end{ex}






\end{document}
